%! From Combinations to Solutions --- Setting Up Ax = b — Unit 2, Lesson 2.0 — Homework (Answer Key)
\documentclass[10pt]{article}
\usepackage{linalg-article}
\usepackage{linalg-key}   % pulls in -boxes; do NOT also load -boxes
\newcommand{\TallMath}[1]{$\displaystyle #1\rule[-1.4em]{0pt}{3.2em}$}
\newcommand{\xx}{\mathbf{x}}
\newcommand{\bb}{\mathbf{b}}

\begin{document}

\pageheader{Unit 2, Lesson 2.0}{Homework --- Answer Key}
\namedateperiod

\vspace{0.12in}

\begin{notesbox}{Practice}
Show your steps and \textbf{check} every solution by rebuilding $\bb$.

\begin{enumerate}[leftmargin=1.9em, itemsep=11pt]
  \item \textbf{Three views.} Let $A=\begin{bmatrix} 2 & 1 \\ 1 & 4 \end{bmatrix}$ and
        $\bb=\begin{bmatrix} 8 \\ 11 \end{bmatrix}$.
    \begin{enumerate}[label=(\alph*), leftmargin=1.8em, itemsep=4pt]
      \item Write $A\xx=\bb$ as a \emph{combination of the columns} of $A$.
            \ans{$x_1\begin{bsmallmatrix} 2\\1 \end{bsmallmatrix}+x_2\begin{bsmallmatrix} 1\\4 \end{bsmallmatrix}=\begin{bsmallmatrix} 8\\11 \end{bsmallmatrix}$}
      \item Write $A\xx=\bb$ as a \emph{system of two equations}.
            \ans{$2x_1+x_2=8$;\ \ $x_1+4x_2=11$}
    \end{enumerate}
  \item \textbf{Solve by elimination.} Find $x_1$ and $x_2$ for each system.
    \begin{enumerate}[label=(\alph*), leftmargin=1.8em, itemsep=4pt]
      \item \ans{$x_1=2,\ x_2=4$} \hfill
      \item \ans{$x_1=4,\ x_2=2$} \hfill
      \item \ans{$x_1=3,\ x_2=2$}
    \end{enumerate}
    {\footnotesize (a) subtract eq1 from eq2: $x_1=2$, then $x_2=4$.\ \ (b) subtract eq2 from
    eq1: $2x_1=8$, $x_1=4$, then $x_2=2$.\ \ (c) add the equations: $3x_1=9$, $x_1=3$, then $x_2=2$.}
  \item \textbf{Fertilizer blend.} Each bag of \textbf{Base A} has $1$ kg nitrogen and $2$ kg
        phosphorus; each bag of \textbf{Base B} has $3$ kg nitrogen and $1$ kg phosphorus. A
        field needs $13$ kg nitrogen and $11$ kg phosphorus.
    \begin{enumerate}[label=(\alph*), leftmargin=1.8em, itemsep=6pt]
      \item Set up and solve $A\xx=\bb$ for the number of bags of each base.
            \ansline{$x_1+3x_2=13,\ 2x_1+x_2=11$. Double eq1: $2x_1+6x_2=26$; subtract eq2:
            $5x_2=15\Rightarrow x_2=3$, then $x_1=4$. Check: $4\begin{bsmallmatrix} 1\\2 \end{bsmallmatrix}+3\begin{bsmallmatrix} 3\\1 \end{bsmallmatrix}=\begin{bsmallmatrix} 13\\11 \end{bsmallmatrix}$ \checkmark}
      \item \emph{Interpret:} how many bags of each, and does the answer make sense?
            \ansline{$4$ bags of Base A and $3$ of Base B --- whole, positive bags, so it is a workable blend.}
    \end{enumerate}
  \item \textbf{How many solutions?} Classify each system as \emph{one}, \emph{none}, or
        \emph{infinitely many}, and say why (parallel lines, same line, or crossing).
    \begin{enumerate}[label=(\alph*), leftmargin=1.8em, itemsep=4pt]
      \item \ans{Infinitely many} \hfill
      \item \ans{None} \hfill
      \item \ans{One: $(1,2)$}
    \end{enumerate}
    {\footnotesize (a) eq2 is exactly $2\times$ eq1 --- same line.\ \ (b) $2\times$ eq1 would give
    $10$, not $7$ --- parallel, no crossing.\ \ (c) subtract eq1 from eq2: $x_2=2$, then $x_1=1$ --- lines cross once.}
\end{enumerate}
\end{notesbox}

\begin{extensionbox}
\textbf{One variable at a time.} Solve this $3\times 3$ system by eliminating $x$ first:
\[
  \begin{aligned}
    x + y + z &= 6 \\
    x + 2y + z &= 8 \\
    2x + y + z &= 7
  \end{aligned}
\]
Subtract the first equation from the second, and subtract twice the first from the third. You
are left with a small system in $y$ and $z$ --- solve it, then back-substitute for $x$.
\ansline{Eq2$-$Eq1 gives $y=2$. Eq3$-2\cdot$Eq1 gives $-y-z=-5$, i.e.\ $y+z=5$, so $z=3$.
Back-substitute: $x=6-2-3=1$. Solution $(x,y,z)=(1,2,3)$; check all three equations.}
\end{extensionbox}

\begin{spiralbox}
\textbf{Coming next --- Lesson 2.1: The Idea of Elimination.} Today you subtracted one equation
from another to make a variable disappear. Next we turn that into a reliable, step-by-step
\emph{procedure} --- choosing \textbf{pivots} and clearing out the entries below them --- so it
works for any size system, not just the lucky ones.
\end{spiralbox}

\end{document}
